% Ad Familiares 10.25 (ad-familiares-10-25)
% Source: https://raw.githubusercontent.com/PerseusDL/canonical-latinLit/master/data/phi0474/phi056/phi0474.phi056.perseus-lat1.xml
% Fetched by scripts/fetch_latin.py

% Dateline (Perseus): Scr. Romae circ. vii K. Iun. a. 711 (43).

\ciceroLetterOpener{CICERO S. D. FVRNIO.}

\ciceroSection{1} si interest, id quod homines arbitrantur, rei p. te, ut instituisti atque fecisti, navare operam rebusque maximis, quae ad exstinguendas reliquias belli pertinent, interesse, nihil videris melius neque laudabilius neque honestius facere posse, istamque operam tuam, navitatem, animum in rem p. celeritati praeturae anteponendam censeo. nolo enim te ignorare quantam laudem consecutus sis ; mihi crede, proximam Planco, idque ipsius Planci testimonio, praeterea fama scientiaque omnium.

\ciceroSection{2} quam ob rem, si quid operis tibi etiam nunc restat, id maximo opere censeo persequendum ; quid enim honestius, aut quid honesto anteponendum? sin autem satis factum rei p. putas, celeriter ad comitia, quoniam mature futura sunt, veniendum censeo, dum modo ne quid haec ambitiosa festinatio aliquid imminuat eius gloriae, quam consecuti sumus. multi clarissimi viri, cum rei p. darent operam, annum petitionis suae non obierunt. quod eo facilius nobis est, quod non est annus hic tibi destinatus, ut, si aedilis fuisses, post biennium tuus annus esset. nunc nihil praetermittere videbere usitati et quasi legitimi temporis ad petendum. video autem Planco cos., etsi etiam sine eo rationes expeditas haberes, tamen splendidiorem petitionem tuam, si modo ista ex sententia confecta essent.

\ciceroSection{3} omnino plura me scribere, cum tuum tantum consilium iudiciumque sit, non ita necesse arbitrabar ; sed tamen sententiam meam tibi ignotam esse nolebam ; cuius est haec summa, ut omnia te metiri dignitate malim quam ambitione maioremque fructum ponere in perpetuitate laudis quam in celeritate praeturae. haec eadem locutus sum domi meae adhibito Quinto, fratre meo et Caecina et Calvisio studiosissimis tui, cum Dardanus, libertus tuus, interesset. omnibus probari videbatur oratio mea ; sed tu optime iudicabis.
